\section{Caso de estudio: Offline RL en el mundo real}

\subsection{Ingeniería compleja de reward}

En los sistemas de interacción humano-máquina, el reward suele incluir múltiples objetivos (accuracy, fairness, latency). El expositor describe un proyecto de LinkedIn: en un sistema de notification se optimizan simultáneamente CTR y retention, por lo que es necesario etiquetar los reward components en el dataset para permitir un entrenamiento descompuesto posterior.

\begin{importantbox}{Truco práctico para reward multiobjetivo}
En el dataset se precalcula un reward vector (CTR, retention, quality), y durante el offline training se usa una linear combination más CVaR regularization para controlar los worst-case outcomes.
\end{importantbox}

\subsection{Demostración visual}

\begin{figure}[H]
\centering
\includegraphics[page=6,width=0.9\textwidth]{07_cs224r_offline_rl_2025.pdf}
\caption{Reality check que muestran las slides: alineación entre el reward multiobjetivo y las métricas de deployment\protect\footnotemark}
\end{figure}
\footnotetext{Página 6 de las slides, que enfatiza la correspondencia uno a uno entre el reward vector y las deployment metrics.}

\subsection{Resumen de la sección}

En los casos reales el reward es multidimensional; es indispensable precisar a nivel de los datos el significado de cada dimensión, o de lo contrario la offline policy quedará desalineada.

